\documentclass[11pt]{article}
\usepackage[T1]{fontenc}
\usepackage{textcomp}
\usepackage[margin=0.75in]{geometry}
\usepackage{amsmath,amssymb,amsthm}
\usepackage{array,booktabs}
\usepackage{hyperref}
\setlength{\parindent}{0pt}
\newtheorem{theorem}{Theorem}
\theoremstyle{definition}
\newtheorem{definition}{Definition}
\title{Flow Proof Artifact: prop-25-describe-circle-from-segment}
\author{Generated by \texttt{flow doc proof}}
\date{Flow Proof Book}
\begin{document}
\maketitle
Given a segment of a circle, to describe the complete circle of which it is a part.\\[0.5em]
\textbf{Source.} euclid — Elements, Book III, Proposition 25\\[0.5em]
\bigskip
\noindent{\large \textbf{Derived fact 1.} \textit{Given a segment of a circle, to describe the complete circle of which it is a part} \hfill the Euclidean plane \cdot Euclid Book III \cdot Proposition 25: to describe the circle of which a given segment is part}\par
\smallskip\noindent\textit{Coordinate: the Euclidean plane \cdot Euclid Book III \cdot Proposition 25: to describe the circle of which a given segment is part}\par
\smallskip\noindent\textit{Source: euclid — Elements, Book III, Proposition 25}\par
\smallskip\noindent\textit{Built on: proposition 21: angles in the same segment are equal, for Euclid Book III on the Euclidean plane, proposition 9: bisect a given angle, for Book I of the Elements in on the Euclidean plane}\par
\medskip
\noindent\textbf{Goal.} Given a segment of a circle, to describe the complete circle of which it is a part\par
\noindent$\displaystyle \text{circle from segment exists}$\par
\medskip
\noindent
\begin{tabular}{@{} >{\raggedright\arraybackslash}p{0.06\textwidth} >{\raggedright\arraybackslash}p{0.47\textwidth} >{\raggedleft\arraybackslash}p{0.41\textwidth} @{}}
\textbf{\#} & \textbf{Proof} & \textbf{Mathematics} \\
\midrule
\textbf{1.} & We prove that proposition 25: to describe the circle of which a given segment is part for Euclid Book III the Euclidean plane. & \\[0.45em]
\textbf{2.} & Let seg = given\_segment. & \\[0.45em]
\textbf{3.} & Let chord = base\_of\_segment. & \\[0.45em]
\textbf{4.} & Let center = point\_equidistant\_from\_ends. & \\[0.45em]
\textbf{5.} & From step 2, step 3, and step 4, we can deduce that circle through segment exists. & $\displaystyle \text{circle through segment exists}$ \\[0.45em]
\textbf{6.} & We invoke the derived fact governing Euclid Book III the Euclidean plane: proposition 21: angles in the same segment are equal, for Euclid Book III on the Euclidean plane. & \\[0.45em]
\textbf{7.} & We invoke the derived fact governing Book I of the Elements in the Euclidean plane: proposition 9: bisect a given angle, for Book I of the Elements in on the Euclidean plane. & \\[0.45em]
\textbf{8.} & From step 6 and step 7, this implies circle from segment exists. Hence proven. & $\displaystyle \text{circle from segment exists}$ \\[0.45em]
\bottomrule
\end{tabular}
\label{thm:Geometry-Euclid Book III-proposition_25_to_describe_the_circle_of_which_a_given_segment_is_part}
\par\medskip\hrule
\end{document}
